\section{Finite dihedral census}

The executable audit realizes the theorem on the (2n) residues for every
(3\leq n\leq32). It takes
\[
H(x)=x+1,
\qquad
A(x)=-x,
\qquad
B(x)=-x-1
\pmod{2n},
\]
so that (A^2=B^2=1) and (AB=H).

For every tested (n), the audit enumerates the generated permutation group
and checks:
\begin{enumerate}
\item the group order is (4n);
\item all (4n) dihedral normal forms are distinct;
\item (BA=H^{-1}) and ([A,B]=H^2);
\item the commutator subgroup has order (n);
\item (	au=H^n) is a central involution;
\item (	au\in\langle H^2\rangle) exactly for even (n);
\item the even and odd reconstruction formulas hold;
\item the Klein-four quotient has kernel (langle H^2\rangle);
\item the quotient image of (	au) is parity-correct.
\end{enumerate}

The census is a reproducible implementation witness. The theorem for every
(n\geq3) is established by the symbolic proofs in the main text.

